\raggedright
\chapterpage{EVIDENCE APPENDIX\enspace /\enspace I}{How to find the arguments}
The keys below identify the native sources. For E1--E6, directory names are relative to \pathpart{research/notes/}. Locators name a proof by identifier or section title. The present export checked source correspondence and typesetting. Only the proof of MC-S1/S2 on page 5 is reproduced in full.

\subhead{E1\enspace /\enspace Three branches and returns}
\pathpart{2026-09-16-three-branch-lift/README.md}, \S1 (maps, $4\to22$, literal iteration); \pathpart{2026-09-16-three-branch-continuation/README.md} (bounded continuation). See \pathpart{2026-09-21-growing-step-cubic/structure-and-reuse.md}, GS-P3, for selected coverage of the output--difference grid. Exact portable counts are in \pathpart{examples/2026-09-16-three-branch-lift.json} from the repository root.

\subhead{E2\enspace /\enspace Sharing geometry}
\pathpart{2026-09-30-elliptic-geometry/genus.md}, EG-G1 and EG-G2, gives the fiber parameter and the sharing classification. \pathpart{complex-extensions.md}, EG-C1/C2 in the same directory, gives the Weierstrass transform and infinite-order point. Its \pathpart{README.md}, sharing-prefix table, records the 23-scalar example. See also the dated \pathpart{verification.md}. \cite{elliptic}

\subhead{E3\enspace /\enspace Gluing and the first frame bridge}
\pathpart{2026-10-06-sharing-compatibility/construction.md}, SH-C1--C3, and \pathpart{2026-10-06-variable-gap-sharing/construction.md}, the band construction and coverage ledger, give the two network counts. \pathpart{2026-10-06-changed-collision-pattern/rigidity.md}, CP-R1, is the fixed-incidence rigidity proof. In the same directory, \pathpart{construction.md}, CP-C1--C4, proves the conic identity, six-scalar bridge and ladder bound. \pathpart{exact-audit.md} records the full-square data.

\subhead{E4\enspace /\enspace Admissible points and arithmetic}
\pathpart{2026-10-11-torsion-compatible-sharing/geometry.md}, TC-G1--G6, separates the finite completed-frame locus from the genus-four cover for a general row. Its \pathpart{exact-results.json} and \pathpart{exact_audit.py} retain the two six-point squares and bounded torsion checks. \pathpart{2026-10-11-marked-shift-certificate/arithmetic.md}, MS-C1, resolves the marked shift at all orders. \pathpart{referee-audit.md} records the historical internal review. The proof uses one specified good place over 23. \cite{frames}

\subhead{E5\enspace /\enspace Varying outputs and excluded closures}
\pathpart{2026-10-11-varying-output-frames/compatibility.md}, VO-C1--C5; \pathpart{2026-10-11-vo-o1-freeze/README.md}, \S3, VO-C6; \pathpart{2026-10-20-tetrahedral-junction/compatibility.md}, TJ-C1--C4, with \pathpart{catalogue.md} for the four roles; \pathpart{2026-10-21-tj-o1-freeze/README.md}, \emph{Complete contradiction}, TJ-C6, with \pathpart{proof-check.md} for the internal check.

\chapterpage{EVIDENCE APPENDIX\enspace /\enspace II}{Status, data and evidence boundaries}
\subhead{E6\enspace /\enspace Exact branching}
\pathpart{research/notes/2026-10-21-multiscale-collisions/}: \pathpart{specialist.md}, MC-S1--S3; \pathpart{README.md}, \S\S1--3; \pathpart{experiment-summary.json} and \pathpart{referee-counts.json}. Six inherited complete-square checks compare 533 ordered-pair outputs. The archival build ran no research grid. The local source for the proposed Raz--Zahl Theorem 1.18 comparison is marked \tagword{MISSING SOURCE}. No exponent from that comparison enters MC-S1/S2. \cite{scales}

\subhead{E7--E10\enspace /\enspace Track-specific sources}
\textbf{E7.} \pathpart{research/tracks/02-quartic-hybrid/STATUS.md}, \emph{Accepted C2--C4 results} and \emph{Carrier and signed rate}, specifies the local ring, counted set and full-interval packing obstruction. The proof is reported from the accepted C4 checkpoint. The status file identifies its local provenance.

\textbf{E8.} \pathpart{research/tracks/03-eisenstein-cartesianization/STATUS.md}, \emph{PROVED adaptive construction and exact finite ledger} and \emph{OPEN selected obligation and CONDITIONAL consequence}. Under \pathpart{research/notes/}, see \pathpart{2026-10-11-bff-o1/README.md}, \S\S1--5; \pathpart{2026-10-21-relative-unit-population/arithmetic.md}, MSU-A1--A3; and \pathpart{2026-10-21-typical-unit-returns/README.md}, equations (11)--(15), with \pathpart{incidences.md} for the orbit partition.

\textbf{E9.} \pathpart{research/tracks/04-cubic-slices/STATUS.md}, \emph{PROVED input theorem; OPEN full output} and \emph{PROVED exact pair evolution; missing output count}. Its named input proof and audit are identified local records. A depth-30 original log remains a \tagword{MISSING PROVENANCE FILE}.

\textbf{E10.} \pathpart{docs/release/2026-11-05/shared-quadratic.md}, Theorem 1, MP-01--MP-07, gives the public mathematical record of the 4 November audit. The original archive manifest and source-conflict register are \pathpart{research/manifest.json} and \pathpart{research/MIGRATION_CONFLICTS.md}. Certain underlying local reference PDFs, administrative records and large raw outputs are retained separately. Their exclusion from a public export is distinct from missing provenance.

\smallskip\tagword{Reading the labels.} PROVED means a complete argument is identified in the source system. CONDITIONAL conclusions require the named hypotheses. FAILED describes only the specified attempt. EXPERIMENTAL counts have finite range. OPEN identifies a statement without resolution in the record. Historical internal checks are attributed at their original dates. This export certifies no external review or novelty.

\renewcommand{\refname}{Selected native reports}
\begingroup\small\setlength{\parskip}{0pt}\bibliographystyle{unsrt}\bibliography{references}\endgroup
